\question
Студент Лёша плохо учил дискретную математику, в итоге попал на допсу и был отчислен. Чтобы зарабатывать на жизнь, он устроился в ресторан «Вкусно – и точка» кассиром. Как-то Лёша заскучал на кассе и решил себя развлечь.

В ресторане есть 3 типа товаров: бургеры, снеки, и напитки. Всего существует $A$ видов бургеров, $B$ видов снеков и $C$ видов напитков. Лёша считает заказ «вкусным», если в нём содержатся блюда хотя бы 2 разных типов, причём для каждого типа в заказе должно быть ровно одно блюдо. Помогите Лёше посчитать, сколько всего существует возможных «вкусных» заказов.

---------------

Автор -- Тимур Гонтарь
